\documentclass[a4paper]{article}
\usepackage{amsfonts}
\usepackage{amsmath}
\usepackage{amssymb}
\usepackage{graphicx}     

\begin{document}
  
\noindent \large Auteur: Shahin Jamal \\
\noindent \large Studierichting: Informatica\\
\noindent \large Oefeningenreeks 1D nr. 19.13\\

\medskip

\normalsize


De nulpunten bepalen van: $z^5 - 3z^4 + 7z^3 - 13z^2 + 12z - 4$\\

Je kan 3 keer na elkaar 1 invullen, dus $(z-1)^3$ is een deler.

\[
\begin{tabular}{c|ccccccccc}
& $1$ & $-3$ & $7$ & & $-13$ & & $12$ & & $-4$ \\
$1$ & & $1$ & $-2$ & & $5$ & & $-8$ & & $4$ \\ \hline
& $1$ & $-2$ & $5$ & & $-8$ & & $4$ & \vline & $0$ \\
$1$ & & $1$ & $-1$ & & $4$ & & $-4$ & & \\ \cline{1-8}
& $1$ & $-1$ & $4$ & & $-4$ & \vline & $0$ & & \\
$1$ & & $1$ & $0$ & & $4$ & & & & \\ \cline{1-7}
& $1$ & $0$ & $4$ & \vline & {$0$} & & & &
\end{tabular}
\]

We krijgen:
\[
(z^2 + 4)(z - 1)^3 = 0
\]

\( z^2 + 4 = 0 \) geeft de oplossingen:
\[
z = \pm 2i
\]

De nulpunten zijn:
\[
\{ 1, 2i, -2i \}
\]

en de ontbinding in factoren is: 
\[
(z^2 + 4)(z - 1)^3
\]

\begin{thebibliography}{999}
%\bibitem{a} Referentie
\end{thebibliography}
\end{document} 
